% !TeX root = ../main.tex
\renewcommand{\abstractname}{Resumen}
\begin{abstract}
En el marco de la transición hacia la Industria 4.0, los activos industriales y vehículos de guiado autónomo han evolucionado hacia sistemas capaces de emitir volúmenes continuos de telemetría en tiempo real, generando un gran caudal de información. Sin embargo, la simple recolección de estos datos no aporta valor operativo por sí sola. Únicamente mediante un procesamiento y análisis adecuados es posible transformar dichos datos en conocimiento útil para la toma de decisiones estratégicas.
De ahí surge la idea de este proyecto, desarrollado en colaboración con EDAG, entidad de referencia en el sector industrial.\\ 

A pesar del avance en el mantenimiento predictivo durante los últimos años, existe una brecha notable entre los estudios teóricos sobre detección de anomalías y su aplicación efectiva en entornos de producción. Esta limitación se debe principalmente a la complejidad en la recolección de datos de calidad, al dinamismo de los entornos industriales y a la necesidad de procesar la información con bajas latencias para actuar antes de que ocurra una avería.\\

Para dar respuesta a esta problemática, este proyecto presenta una solución integral para la detección de anomalías en tiempo real sobre una flota de robots móviles. El sistema se compone de una infraestructura distribuida de ingesta y procesamiento continuo de datos, una capa de almacenamiento histórico optimizada para series temporales y un módulo analítico basado en aprendizaje automático para la identificación de comportamientos anómalos. Todo el conjunto se completa con una herramienta visual de monitorización y alertas para el personal de planta, desarrollada bajo una arquitectura limpia, desacoplada y escalable.

\medskip
\noindent\textbf{Palabras clave:} IIoT, Mantenimiento Predictivo, Detección de Anomalías, Aprendizaje Automático, Robots Móviles, Industria 4.0, Data Lakehouse, Arquitectura Medallón.
\end{abstract}